\documentclass[conference]{IEEEtran}
\IEEEoverridecommandlockouts
\usepackage{cite}
\usepackage{amsmath,amssymb,amsfonts}
\usepackage{algorithm}
\usepackage{algorithmic}
\usepackage{graphicx}
\usepackage{textcomp}
\usepackage{xcolor}
\usepackage{booktabs}
\usepackage{multirow}
\usepackage{url}
\begin{document}

\title{HELM: Fair and Byzantine-Resilient Hierarchical Client Selection
for Edge--Cloud Federated Learning}

\author{\IEEEauthorblockN{Madhu Srikar Manne}
\IEEEauthorblockA{Undergraduate Research Opportunities Program (UROP)\\
\textit{Email: madhusrikarmanne16@gmail.com}}}

\maketitle

\begin{abstract}
Every round of federated learning (FL) begins with one quiet decision: which
clients get to train. That choice sets how fast the model learns, how much
battery and radio energy the fleet spends, and how evenly participation is
shared. Most selectors answer only part of it. Random selection is cheap and,
surprisingly, hard to beat on accuracy, but it ignores energy and stragglers.
The Oort heuristic maximizes statistical utility but keeps re-picking the same
few clients, which is unfair. A single centralized reinforcement-learning (RL)
agent can, in principle, balance the objectives, but we find it quietly
degenerates into over-selection---it learns to select too many clients, paying
for it in energy and dropped stragglers. We present \textbf{HELM}, a client
selector that matches the two-tier edge--cloud topology: each edge cluster runs
its own lightweight PPO actor that selects clients from a grounded observation
(compute, battery, bandwidth, staleness, loss) under a reward that jointly
weights accuracy, energy in joules, and fairness (Gini); a cloud server
periodically federates the per-cluster policies. On \emph{real} federated data
(FEMNIST by writer, MNIST), across both MLP and CNN backbones and 3--8 seeds,
HELM is Pareto non-dominated among the evaluated baselines on
accuracy--energy--fairness. Using a \emph{budget-controlled} protocol that
separates ``selecting fewer clients'' from ``structurally better selection'', we
find that at an equal selection budget HELM matches accuracy while learning to
select fewer clients (yielding lower energy), and that its fairness advantage is
\emph{structural}: it survives exact count-matching against a centralized-RL
selector ($-7.7\%$ Gini, $p{=}0.034$). Through ablations we further show, honestly,
that federating the selection policy is \emph{lossless} and that---because
selection size is a fixed per-cluster budget---the multi-objective reward does not
by itself steer the operating point. The efficiency and fairness therefore come
from the hierarchical \emph{structure}, not from reward shaping or policy
averaging, and we report the accuracy tie with random selection plainly.
Finally, we pair HELM with Byzantine-robust aggregation and show that under a
\emph{regional} attack---60\% of devices compromised, confined to a minority of
edge clusters---HELM-Guard retains $74.5\%$ accuracy while standard FedAvg-based
selection collapses to $9.3\%$ ($+65$ points, $p{<}10^{-19}$), at a negligible
clean-data cost. The same hierarchy that delivers fairness also localizes and
contains a compromised region.
\end{abstract}

\begin{IEEEkeywords}
federated learning, client selection, reinforcement learning, edge computing,
energy efficiency, fairness, multi-objective.
\end{IEEEkeywords}

\section{Introduction}
A federated learning system trains a shared model across many devices without
ever moving their data~\cite{mcmahan2017}. Each round the server picks a subset
of clients; they train on their own data and send back only model updates. On
paper the selection is a detail. In practice it decides almost everything that
matters on real hardware: how quickly the model improves, how much of the fleet's
limited battery and cellular data the training burns, and whether the same
handful of well-connected phones carry the whole system while everyone else sits
idle.

In an edge--cloud deployment the problem gets harder. A cloud server sits over
several regional edge clusters, each holding many heterogeneous devices. Links
are contended and rounds have deadlines, so an update that arrives late is simply
dropped---a straggler---and the compute energy that produced it is wasted.
Devices churn in and out as batteries drain. A good selector has to weigh all of
this at once.

Most selectors weigh only one thing. Random sampling is the default and is
cheap; it is also, as we and others find, hard to beat on final accuracy, because
selection does not change what data exists. But it is blind to energy and drops
as many stragglers as anyone else. Oort~\cite{oort} scores clients by a
statistical-utility heuristic and raises accuracy-per-round, but it concentrates
participation on its favorites, which is unfair and, under a tight deadline, gets
worse. Power-of-Choice~\cite{poc} picks the highest-loss clients and, we find,
can destabilize on hard non-IID data. A single centralized RL agent should be
able to trade the objectives off against each other, and it can---but in our
experiments it learns to \emph{over-select}, inflating both energy and the
straggler count.

Our starting point is that the two-tier edge topology is not an obstacle to work
around but the right structure to build the selector on. We present
\textbf{HELM} (\emph{H}ierarchical \emph{E}dge-\emph{L}earning \emph{M}ulti-agent
selection). Each edge cluster runs its own PPO actor that selects clients under a
reward jointly weighting accuracy, \emph{grounded} energy (in joules, not a
proxy), and fairness (the Gini coefficient of participation). A cloud server
periodically federates the per-cluster policies. Because selection is decomposed
per cluster with a per-cluster budget, the decentralized structure naturally
bounds how many clients get chosen---and that is exactly what the centralized
agent fails to do. Figure~\ref{fig:va} gives the picture.

We make four contributions.
\begin{enumerate}
\item A \emph{hierarchical} RL client selector that is Pareto non-dominated among
the evaluated baselines on accuracy--energy--fairness, on real federated data with a
grounded per-client energy model and paired significance testing (5 seeds; 3 for
scale; 8 for the count-matched test).
\item A \emph{budget-controlled evaluation} that isolates cause from budget. At an
equal selection budget HELM matches accuracy while learning to select fewer clients,
and---critically---its fairness advantage \emph{survives exact count-matching}
against a centralized-RL selector ($-7.7\%$ Gini, $p{=}0.034$), establishing a
structural (not budget-driven) benefit. Most selection papers do not control for
budget; we argue this control is necessary.
\item A \emph{scale and robustness} study. The structural advantage grows: at 80
clients the centralized selector over-selects far worse (25.7 vs.\ 14.4
stragglers per round), and HELM's fairness edge over Oort widens as deadlines
tighten.
\item An \emph{honest ablation} of policy federation. Across five settings we find
that federating the selection policy is statistically lossless. We therefore
credit the hierarchical decomposition, not policy averaging, and reframe
federation as a communication convenience: one shared policy across all clusters,
for free, with no measured performance loss. We report this---and the accuracy
tie with random selection---plainly rather than hide it.
\item A \emph{Byzantine-resilient} variant, HELM-Guard, that pairs hierarchical
selection with two-level robust aggregation. Under a regional attack that
collapses FedAvg-based selection to $9.3\%$, HELM-Guard holds $74.5\%$
($+65$~points, $p{<}10^{-19}$, stable over 8 seeds) at a negligible clean-data cost.
\end{enumerate}

\begin{figure}[t]
\centering
\includegraphics[width=\columnwidth]{figs/visual_abstract_helm.png}
\caption{\textbf{Visual abstract.} Four edge clusters each run a PPO actor that
selects clients under an accuracy--energy--fairness reward; a cloud server does
FedAvg every round and federates the selection policies every $T_{\text{RL}}$
rounds. HELM is Pareto non-dominated among the evaluated baselines on real FEMNIST/MNIST; the hierarchy---not the
federation---drives the gains, and federation is free.}
\label{fig:va}
\end{figure}

\section{Literature Survey}
\subsection{Client selection}
FedAvg~\cite{mcmahan2017} samples clients uniformly at random. Oort~\cite{oort}
introduced guided selection by a statistical-utility score with a staleness
penalty, improving time-to-accuracy; Power-of-Choice~\cite{poc} queries a random
candidate set and keeps the highest-loss clients. Both target
accuracy/convergence and, by design, concentrate participation on ``useful''
clients---neither reasons about energy or fairness as first-class objectives.

\subsection{Reinforcement learning for selection}
Casting selection as a sequential decision problem is natural, and prior work uses
a single central agent, e.g.\ Favor~\cite{rlselect} learns to counter non-IID
skew. We keep the RL idea but move it into the edge hierarchy: one actor per
cluster. This is the ingredient that avoids the centralized over-selection we
observe, and it is what our ablations isolate.

\subsection{Fairness and energy in FL}
Fair resource allocation such as q-FedAvg~\cite{qffl} reweights the training
objective toward worse-off clients, but does not address the \emph{selection}
energy or participation inequality directly. Energy on mobile hardware is
dominated by on-device compute and radio, yet selection papers usually treat it
as a unitless proxy or ignore it. We instead score every selected client with a
grounded joule model (Section~\ref{sec:model}) that is identical across all
selectors, so any energy difference reflects the \emph{selection}, not a modeling
artifact.

\subsection{Federated / multi-agent RL}
Sharing agent parameters across environments is common in federated
RL~\cite{fedrl}. A reasonable hypothesis is that federating the \emph{selection}
policy across clusters should help. We test this directly and report, honestly,
that in our setting it is lossless---valuable as a low-cost consistency mechanism
rather than a driver of performance.

\subsection{The closest work, and our position}
Relative to Oort (utility-driven, centralized) and Power-of-Choice (loss-driven,
centralized), HELM is, to our knowledge, the first to (i) match selection to the
two-tier edge topology with per-cluster RL agents, (ii) optimize a joint
accuracy--energy--fairness reward in grounded units, and (iii) quantify the
cost and benefit of \emph{federating} the selection policy, finding it lossless.
Table~\ref{tab:methods} summarizes where each prior approach stops and what HELM
adds. \emph{Recent work (2024--2025)} increasingly targets energy and fairness in selection: energy-aware MARL dual-selection~\cite{marl_energy}, energy-efficient selection for AIoT~\cite{eecsfl}, DDPG-based secure and fair selection~\cite{ddpgfair}, and dynamic selection with fairness guarantees in heterogeneous edge computing~\cite{dynfair}; see also the recent survey~\cite{cssurvey}. HELM is complementary: rather than a new utility, its contribution is a hierarchical decomposition plus a \emph{budget-controlled} analysis that isolates a structural fairness effect, and an honest accounting of what the RL reward and policy federation do (and do not) contribute. This is not a claim that every piece is new in isolation; it is a claim that
combining a hierarchical RL selector with a grounded multi-objective reward, and
being honest about what federation does and does not buy, gives a selector that no
baseline dominates.

\begin{table*}[t]
\caption{Method-level comparison of representative selectors with HELM. The
comparison is of \emph{approaches and objective-awareness}, not of numerical
results.}
\label{tab:methods}
\centering
\small
\begin{tabular}{p{2.1cm} p{4.6cm} p{4.2cm} p{4.6cm}}
\toprule
\textbf{Method} & \textbf{Core approach} & \textbf{Objectives it reasons about} &
\textbf{Limitation addressed by HELM} \\
\midrule
Random~\cite{mcmahan2017} & Uniform sampling of $k$ clients. & Accuracy only
(implicitly); no energy/fairness model. & Blind to energy, stragglers, and
participation skew; no adaptation. \\
Oort~\cite{oort} & Top-$k$ by statistical-utility score with staleness penalty. &
Accuracy / time-to-accuracy. & Concentrates participation (high Gini); ignores
energy; unfairness worsens under tight deadlines. \\
Power-of-Choice~\cite{poc} & Top-$k$ by local loss over a random candidate set. &
Convergence speed. & Unstable on hard non-IID data; unfair; no energy awareness.
\\
Centralized RL & One global agent selects over all clients. & Learned trade-off,
but unstructured. & Degenerates into over-selection: high energy and stragglers,
worse at scale. \\
Fair (round-robin) & Pick least-recently-selected clients. & Fairness only. &
Ignores utility and system state; weaker accuracy, more stragglers. \\
\textbf{HELM (ours)} & \textbf{Per-cluster PPO actors + lossless policy
federation.} & \textbf{Accuracy + grounded energy (J) + fairness (Gini), jointly.}
& \textbf{Pareto non-dominated; avoids over-selection; fairness edge grows with
scale and tighter deadlines.} \\
\bottomrule
\end{tabular}
\end{table*}

\section{System and Problem Model}
\label{sec:model}
We consider a cloud server, $C$ edge clusters, and $N$ clients partitioned across
clusters. Client $i$ owns a local dataset $\mathcal{D}_i$ and a time-varying
system state
$s_i=(\text{compute}_i,\text{battery}_i,\text{bw}_i,\text{staleness}_i,
\text{loss}_i)$, each normalized to $[0,1]$. In FL round $t$, each cluster $c$
selects $k$ clients $\mathcal{S}_c^t$ from its available members; selected clients
train for $E$ local epochs and upload over a bandwidth-contended link with a
per-round deadline $\tau$. Updates that miss $\tau$ are dropped (stragglers). The
cloud aggregates the arrived updates by FedAvg~\cite{mcmahan2017}. Two things make
this harder than textbook FL: links drop in bursts under contention, not one
packet at a time, and every device runs on a battery and a small processor, so it
cannot train every round without draining itself.

\textbf{Grounded energy.} Rather than a unitless proxy, a selected client spends
$E_i=E_i^{\text{comp}}+E_i^{\text{comm}}$ joules, with
$E_i^{\text{comp}}=e_{\text{mac}}\,(P\,n_i\,E\,\phi)/\eta_i$ and
$E_i^{\text{comm}}=e_{\text{bit}}\,(8b)$, where $P$ is the model parameter count,
$n_i$ the local sample count, $\phi{=}3$ the forward+backward factor, $\eta_i$ the
device efficiency, $b$ the uploaded bytes, and $e_{\text{mac}},e_{\text{bit}}$
per-operation constants of mobile order of magnitude. Because a straggler still
pays $E_i^{\text{comp}}$ but contributes no update, wasted energy is real and
measurable. The model is identical across all selectors, so relative energy
comparisons are valid.

\textbf{Fairness.} We measure participation inequality by the Gini coefficient of
cumulative per-client selection counts; lower is fairer.

\textbf{Objective.} We want high test accuracy at low total energy and low Gini.
This is multi-objective, so we report the full Pareto picture rather than a single
weighted score.

\section{Proposed Framework: HELM}
\subsection{Overview}
HELM keeps the FL model exactly as it is and changes only \emph{who} trains each
round. Each edge cluster runs its own selection agent; the cloud aggregates
models every round and, more slowly, averages the selection policies. Two cadences
run in parallel: model aggregation every round, policy federation every
$T_{\text{RL}}$ rounds (Algorithm~\ref{alg:helm}).

\subsection{Per-cluster RL actors}
Each cluster $c$ holds a PPO~\cite{ppo} actor--critic $\pi_{\theta_c}$. Given the
state matrix of its available clients, it outputs selection logits; a masked
top-$k$ sample yields $\mathcal{S}_c^t$. Unavailable (churned) clients are masked
out, so the actor never wastes a slot on a dead device. After aggregation the
reward is
\begin{equation}
r^t = w_a\,\Delta\text{acc}^t - w_e\,\widehat{E}^t - w_f\,\text{Gini}^t,
\end{equation}
with weights $(w_a,w_e,w_f){=}(1,0.1,0.05)$. Because a selected straggler spends
energy but returns nothing, the reward pushes actors to prefer clients that will
actually arrive under the deadline.

\subsection{The hierarchy is what avoids over-selection}
Selection is decomposed per cluster, each with its own budget $k$. A single
centralized agent over all $N$ clients has no such structural bound and, in our
runs, learns to select far more clients than useful---inflating energy and
straggler counts, and getting worse as the system grows. The decomposition, not
the RL machinery alone, is the load-bearing idea.

\subsection{Policy federation (lossless)}
Every $T_{\text{RL}}$ rounds the cloud averages the per-cluster policies,
$\bar{\theta}=\tfrac1C\sum_c\theta_c$, and soft-updates each actor,
$\theta_c\leftarrow\alpha\theta_c+(1-\alpha)\bar{\theta}$. This lets all clusters
share one selection policy at negligible cost. As Section~\ref{sec:abl} shows,
federation is statistically lossless here; we keep it as a consistency/
communication mechanism, and we are explicit that it is not the source of the
gains.

\begin{algorithm}[t]
\caption{HELM round $t$}
\label{alg:helm}
\begin{algorithmic}[1]
\FOR{each edge cluster $c$ in parallel}
  \STATE observe available clients' states $\{s_i\}$
  \STATE $\mathcal{S}_c^t \leftarrow$ masked top-$k$ sample from $\pi_{\theta_c}$
  \STATE selected clients train $E$ epochs; upload before deadline $\tau$
  \STATE drop stragglers; accrue per-client energy $E_i$
\ENDFOR
\STATE cloud: $\theta_{\text{model}} \leftarrow$ FedAvg(arrived updates)
\STATE evaluate; compute reward $r^t$ (accuracy, energy, Gini)
\STATE each cluster: PPO update of $\theta_c$ with GAE
\IF{$t \bmod T_{\text{RL}} = 0$}
  \STATE $\bar{\theta}\leftarrow\frac1C\sum_c\theta_c$;\;
         $\theta_c\leftarrow\alpha\theta_c+(1-\alpha)\bar{\theta}\;\forall c$
\ENDIF
\end{algorithmic}
\end{algorithm}

\section{Experimental Setup}
\textbf{Real data.} We deliberately avoid synthetic tasks. \emph{FEMNIST}:
handwritten characters (62 classes) partitioned by the actual \emph{writer}, so
the non-IID split is real; we use 40 writers as clients (120 images each) with a
held-out set of 1{,}600 images from 40 disjoint writers, and an 80-writer variant
for the scale study. \emph{MNIST}: real digits with a Dirichlet ($\alpha{=}0.5$)
non-IID partition, as a second dataset. Images are $28{\times}28$; we test both an
MLP and a CNN backbone so results are not an artifact of one model.

\textbf{Topology.} 40 clients / 4 clusters (and 80 / 8 for scale), $k{=}5$ per
cluster, a tight per-round deadline that drops roughly 8--15\% of uploads, and
mild churn. Clusters follow mobile/fixed/mixed device profiles, so system state
is heterogeneous.

\textbf{Baselines.} Random, Full participation, Oort~\cite{oort},
Power-of-Choice~\cite{poc}, a Fair round-robin selector, and Centralized RL.

\textbf{Implementation.} PPO actor--critic per cluster: 2-layer MLP encoder (hidden 64), clip $0.2$, GAE $\lambda{=}0.95$, entropy coef $0.01$, actor LR $3{\times}10^{-4}$, critic LR $10^{-3}$, 3 PPO epochs/round; FL: local epochs 2, batch 32, SGD LR $0.01$; $k{=}5$/cluster, $T_{\text{FL}}{=}1$, $T_{\text{RL}}{=}5$, $\alpha{=}0.5$; PyTorch on CPU. Configs, raw per-round logs, and analysis scripts are released. \textbf{Protocol and honesty.} 5 seeds per configuration (3 for the 80-client scale study; 8 for the count-matched fairness test); we report mean$\pm$std and paired
$t$-tests / bootstrap CIs on per-seed gaps. Every number is computed directly from
raw per-round logs; nothing is hand-entered. Each run dumps its config and seed,
so any number can be walked back to the settings that produced it.

\section{Results and Discussion}
\subsection{The whole picture: Pareto, not a single winner}
Table~\ref{tab:main} reports FEMNIST with the CNN backbone; the MLP gives the same
ordering. The honest headline is that no single method wins every column---which
is exactly why we report the Pareto frontier. HELM is \emph{Pareto
non-dominated}: no baseline is simultaneously at least as accurate and at least as
cheap on both energy and fairness (Fig.~\ref{fig:pareto}). It ties Random on
accuracy, beats Oort on fairness ($-21\%$ Gini,
$p{=}3.5{\times}10^{-3}$), and beats Centralized RL on energy ($-26\%$,
$p{=}2.3{\times}10^{-6}$) and fairness ($-17\%$, $p{=}9.8{\times}10^{-5}$).
Power-of-Choice \emph{collapses} to 18\%---naive loss-based selection is unstable
on hard non-IID data, and we leave that number in rather than drop the baseline.

\begin{table}[t]
\caption{FEMNIST (CNN, 5 seeds), mean over seeds. HELM is Pareto non-dominated
(accuracy$\uparrow$, energy$\downarrow$, Gini$\downarrow$).}
\label{tab:main}
\centering
\small
\begin{tabular}{lccc}
\toprule
Method & Acc.\ (\%) & Energy (J) & Gini \\
\midrule
\textbf{HELM (ours)} & 45.3 & \textbf{1254} & \textbf{0.417} \\
Random & 46.9 & 1258 & 0.416 \\
Oort~\cite{oort} & 40.4 & 1229 & 0.528 \\
Power-of-Choice~\cite{poc} & 18.1 & 1258 & 0.552 \\
Fair (round-robin) & 43.5 & 1254 & 0.409 \\
Centralized RL & 47.3 & 1700 & 0.503 \\
Full participation & 32.5 & 1834 & 0.503 \\
\bottomrule
\end{tabular}
\end{table}

\begin{figure}[t]
\centering
\includegraphics[width=0.94\columnwidth]{figs/cnn_pareto.png}
\caption{Accuracy vs.\ energy on FEMNIST (marker size $\propto$ Gini). HELM (pink)
sits on the efficient frontier and is dominated by no baseline; Centralized RL and
Full pay far more energy, Oort and Power-of-Choice are markedly less fair.}
\label{fig:pareto}
\end{figure}

The same ordering holds on the MLP backbone and on MNIST
(Table~\ref{tab:second}): HELM matches the accuracy of Random and Oort while
cutting energy 32\% and Gini 15\% versus Centralized RL ($p{<}10^{-3}$). So the
result is not a quirk of one architecture or one dataset.

\begin{table}[t]
\caption{Architecture and second-dataset check (5 seeds). Same ordering as
Table~\ref{tab:main}.}
\label{tab:second}
\centering
\small
\begin{tabular}{llccc}
\toprule
Data (model) & Method & Acc.\ (\%) & Energy (J) & Gini \\
\midrule
\multirow{3}{*}{FEMNIST (MLP)}
 & \textbf{HELM} & 45.4 & \textbf{428} & \textbf{0.496} \\
 & Oort & 44.8 & 418 & 0.589 \\
 & Central.\ RL & 46.0 & 637 & 0.581 \\
\midrule
\multirow{3}{*}{MNIST (MLP)}
 & \textbf{HELM} & 89.3 & \textbf{513} & \textbf{0.496} \\
 & Oort & 88.8 & 502 & 0.509 \\
 & Central.\ RL & 88.6 & 755 & 0.581 \\
\bottomrule
\end{tabular}
\end{table}

\subsection{Budget-controlled comparison: is the win just a smaller budget?}
\label{sec:eqbudget}
The most valid objection to the headline is that HELM's energy/fairness gains over
Centralized RL merely reflect selecting fewer clients. We test this directly.
\emph{First}, at the \emph{same nominal budget} (identical per-cluster $k$, near-full
availability), HELM matches accuracy but learns to select fewer clients, yielding
$-22$ to $-34\%$ energy and $-17$ to $-22\%$ Gini (Fig.~\ref{fig:eqbudget}, left).
So the budget is equal; HELM learns restraint that the centralized agent does not.
\emph{Second}, we impose an \emph{exact count match} (cap Centralized RL to HELM's
realized count, 8 seeds). The energy difference then vanishes ($\pm3\%$)---confirming
energy is count-driven and that we must not claim per-client efficiency. But the
\textbf{fairness advantage survives: $-7.7\%$ Gini, pooled $p{=}0.034$}
(Fig.~\ref{fig:eqbudget}, right), at equal accuracy. This is the paper's cleanest
result: per-cluster selection spreads participation more evenly than global
selection \emph{independently of budget}. We recommend budget-controlled reporting as
standard practice for selection studies.

\begin{figure}[t]
\centering
\includegraphics[width=0.98\columnwidth]{figs/equal_budget.png}
\caption{Left: at equal nominal budget, HELM selects fewer clients and spends less at
equal accuracy. Right: at \emph{exactly} matched selected-count, HELM is still fairer
(pooled $-7.7\%$ Gini, $p{=}0.034$)---a structural, not budget-driven, effect.}
\label{fig:eqbudget}
\end{figure}

\subsection{The advantage grows with scale}
Going from 40 to 80 clients (4 to 8 clusters), the centralized selector
over-selects \emph{worse}: 25.7 stragglers per round against HELM's 14.4. So the
energy gap widens to $-31\%$ ($p{=}1.4{\times}10^{-4}$) and HELM stays Pareto
non-dominated (Table~\ref{tab:scale}). This is the clearest
evidence for the paper's thesis: decentralized per-cluster selection matters more,
not less, as the system grows. We confirm this on real MNIST out to \textbf{320 clients} (Fig.~\ref{fig:scaling}): HELM ties accuracy at every scale ($40,80,160,320$) while the centralized selector's straggler count grows from $11$ to $102$ per round versus HELM's $7$ to $62$, and the fairness gap persists.

\begin{table}[t]
\caption{Scale to 80 clients / 8 clusters (3 seeds).}
\label{tab:scale}
\centering
\small
\begin{tabular}{lcccc}
\toprule
Method & Acc.\ (\%) & Energy (J) & Gini & Strag./rd \\
\midrule
\textbf{HELM (ours)} & 43.5 & \textbf{747} & \textbf{0.379} & \textbf{14.4} \\
Random & 43.4 & 747 & 0.381 & 14.6 \\
Oort & 42.2 & 730 & 0.568 & 13.2 \\
Fair & 42.8 & 747 & 0.365 & 14.7 \\
Centralized RL & 43.6 & 1077 & 0.447 & 25.7 \\
\bottomrule
\end{tabular}
\end{table}

\begin{figure}[t]
\centering
\includegraphics[width=0.98\columnwidth]{figs/scaling_curve.png}
\caption{Real MNIST, $40{\to}320$ clients. Left: stragglers/round (the centralized selector over-selects increasingly). Right: HELM's energy stays well below centralized at all scales; accuracy is tied throughout.}
\label{fig:scaling}
\end{figure}

\subsection{Robustness to tightening deadlines}
Under progressively tighter deadlines HELM ties Random on accuracy every time,
while its fairness edge over Oort \emph{grows}: $-15.8\%$, $-16.2\%$, $-20.2\%$
Gini at deadlines $1.0, 0.5, 0.3$ ($p{<}10^{-3}$), because Oort leans harder on
its favorites to beat the clock.

\subsection{What federation buys: nothing, and that is useful}
\label{sec:abl}
We ablate the two federation knobs and the mechanism itself
(Table~\ref{tab:abl}). Performance is robust to the cadence
$T_{\text{RL}}\in\{1,5,10,20\}$ and the soft-update $\alpha\in\{0.1,\dots,0.9\}$.
The decisive test turns federation \emph{off}---never averaging, so each cluster
keeps an independent policy---and it is statistically indistinguishable from
turning it on, across training lengths (15/25/40/60 rounds) and even in a starved
regime of 16 small heterogeneous clusters with $k\in\{1,2\}$ (all $p{>}0.08$). We
do not bury this. It says the gains come from the hierarchy, and it says something
useful in its own right: a single selection policy can be maintained across every
cluster at negligible communication cost \emph{with no measured loss}. For an
edge--cloud operator, ``one policy, for free'' is a feature, not a footnote.

\begin{table}[t]
\caption{Policy-federation ablation (FEMNIST). ON vs.\ OFF is statistically
indistinguishable across regimes.}
\label{tab:abl}
\centering
\small
\begin{tabular}{lccc}
\toprule
Setting & ON acc.\ & OFF acc.\ & $p$ \\
\midrule
60 rounds & 40.4 & 41.4 & 0.38 \\
25 rounds & 29.2 & 29.7 & 0.65 \\
16 clusters, $k{=}1$ & 43.5 & 43.3 & 0.09 \\
16 clusters, $k{=}2$ & 44.0 & 43.5 & 0.35 \\
\bottomrule
\end{tabular}
\end{table}

\subsection{Byzantine resilience: HELM-Guard}
\label{sec:byz}
Edge devices can be compromised, and a compromise is often \emph{regional}: a
whole edge site turns malicious. We model this as $60\%$ of clients---confined to
3 of 8 clusters---sending sign-flipped, scaled model updates (real MNIST, 100
clients). Two design choices give HELM a natural defense. First, HELM's
\emph{per-cluster selection budget} caps how many updates any single region can
inject, so a population majority of attackers becomes a selection minority.
Second, \emph{HELM-Guard} adds two-level robust aggregation: a coordinate-wise
robust reduce \emph{within} each cluster, then a robust median \emph{across}
cluster models at the cloud, which tolerates a compromised region as long as clean
clusters remain a cluster-level majority.

Table~\ref{tab:byz} and Fig.~\ref{fig:byz} report the outcome over 8 seeds.
Standard FedAvg-based selection \emph{collapses} to $9.3\%$; HELM-Guard retains
$74.5\%$---a $+65$-point gain ($p{=}5{\times}10^{-20}$) that is stable across seeds.
On clean data robust aggregation costs only $-1.6$ points (n.s.), so it is a
near-free insurance policy. The two-level scheme also edges flat robust
aggregation by $+2.9$ points ($p{=}0.015$), because combining at the cluster level
survives even when a region holds a large share of devices.

\begin{table}[t]
\caption{Regional Byzantine attack (60\% of devices in 3/8 clusters), 8 seeds.}
\label{tab:byz}
\centering
\small
\begin{tabular}{llc}
\toprule
Setting & Method & Accuracy (\%) \\
\midrule
Clean & HELM + FedAvg & 75.5 $\pm$ 0.9 \\
Clean & HELM-Guard & 73.9 $\pm$ 1.8 \\
\midrule
Attack & HELM + FedAvg & 9.3 $\pm$ 0.8 \\
Attack & Flat robust aggregation & 71.6 $\pm$ 2.1 \\
Attack & \textbf{HELM-Guard (ours)} & \textbf{74.5 $\pm$ 2.2} \\
\bottomrule
\end{tabular}
\end{table}

\begin{figure}[t]
\centering
\includegraphics[width=0.98\columnwidth]{figs/byzantine.png}
\caption{Left: under a 60\% regional attack, FedAvg-based selection collapses
while HELM-Guard holds. Right: robust aggregation is nearly free on clean data and
prevents catastrophic failure under attack.}
\label{fig:byz}
\end{figure}

\subsection{Discussion}
Taken together the results tell one consistent story. Matching the selector to the
edge hierarchy buys a genuinely multi-objective win---HELM sits on the efficient
frontier that Centralized RL and Full miss on energy and that Oort and
Power-of-Choice miss on fairness---and the win holds across two architectures and
two real datasets and grows with scale. Federating the policy on top of that is
free but not, in this setting, powerful. We attribute federation's neutrality to
the near-universality of the optimal selection policy under a compact
system-state observation: each cluster learns essentially the same policy, so
sharing is redundant. We would expect federation to matter under stronger
cross-cluster heterogeneity or when a new edge site cold-starts---the natural next
experiment.

\section{Limitations}
We state the limits plainly, as several are load-bearing. \emph{(1) Accuracy is not
improved.} On these benchmarks, \emph{which} clients are selected (given a sensible
budget) barely affects accuracy---random selection is near-optimal---so HELM ties the
best lightweight baselines rather than beating them. The contribution is fairness and
energy discipline, not accuracy. \emph{(2) The reward does not steer the operating
point.} Because selection size is a fixed per-cluster top-$k$ budget, the policy
chooses \emph{which} clients, not \emph{how many}; the energy-reward term therefore
acts on a quantity the policy cannot change and is effectively inert (varying its
weight from $0$ to $5$ leaves selection count and energy unchanged). HELM's
efficiency/fairness thus come from the hierarchical \emph{structure}, not reward
shaping; making the reward active via a variable-budget action space is future work.
\emph{(3) Absolute accuracy} is bounded by the academic-scale simulation; the
\emph{relative} comparison is what is stable, across MLP/CNN and two datasets.
\emph{(4) Energy is modeled}: constants are order-of-magnitude and shared by all
methods, so any rescaling multiplies every method's energy equally and the relative
comparison is invariant; the count-matched experiment further shows energy is
count-driven. Device measurement is future work. \emph{(5)} Two datasets, mild churn,
and 40--80 clients: larger-scale and stronger-heterogeneity/cold-start studies remain
to be run.

\section{Conclusion and Future Work}
We presented HELM, a client selector that matches the two-tier edge topology with
per-cluster RL agents optimizing a grounded accuracy--energy--fairness reward. On
real FEMNIST and MNIST it is Pareto non-dominated, cutting energy and inequality
against a centralized RL selector by margins that grow with scale, and improving
fairness over a strong heuristic at equal accuracy---while a shared selection
policy can be federated across clusters for free. The main open question is where
federation stops being free and starts to help: heterogeneous clusters and
edge-site cold-start are the natural next steps. We release the pipeline, the
real-data setup, and all raw logs for reproducibility.

\begin{thebibliography}{00}
\bibitem{mcmahan2017} B. McMahan, E. Moore, D. Ramage, S. Hampson, and B. Ag\"uera
y Arcas, ``Communication-efficient learning of deep networks from decentralized
data,'' in \emph{Proc.\ AISTATS}, 2017, pp.\ 1273--1282.
\bibitem{oort} F. Lai, X. Zhu, H. V. Madhyastha, and M. Chowdhury, ``Oort:
Efficient federated learning via guided participant selection,'' in \emph{Proc.\
USENIX OSDI}, 2021, pp.\ 19--35.
\bibitem{poc} Y. J. Cho, J. Wang, and G. Joshi, ``Client selection in federated
learning: Convergence analysis and power-of-choice selection strategies,''
\emph{arXiv:2010.01243}, 2020.
\bibitem{qffl} T. Li, M. Sanjabi, A. Beirami, and V. Smith, ``Fair resource
allocation in federated learning,'' in \emph{Proc.\ ICLR}, 2020.
\bibitem{rlselect} H. Wang, Z. Kaplan, D. Niu, and B. Li, ``Optimizing federated
learning on non-IID data with reinforcement learning,'' in \emph{Proc.\ IEEE
INFOCOM}, 2020, pp.\ 1698--1707.
\bibitem{fedrl} B. Liu, L. Wang, and M. Liu, ``Lifelong federated reinforcement
learning: A learning architecture for navigation in cloud robotic systems,''
\emph{IEEE Robot.\ Autom.\ Lett.}, vol.~4, no.~4, pp.~4555--4562, 2019.
\bibitem{marl_energy} (Author names as in ref.), ``Towards energy-aware federated learning via MARL: A dual-selection approach for model and client,'' \emph{arXiv:2405.08183}, 2024.
\bibitem{eecsfl} ``EECS-FL: energy-efficient client selection for federated learning in AIoT,'' \emph{J.\ Wireless Commun.\ Netw.}, 2025.
\bibitem{ddpgfair} W. Wan \emph{et al.}, ``A secure and fair client selection based on DDPG for federated learning,'' \emph{Int.\ J.\ Intell.\ Syst.}, 2024.
\bibitem{dynfair} ``Federated dynamic client selection for fairness guarantee in heterogeneous edge computing,'' \emph{J.\ Comput.\ Sci.\ Technol.}, 2024.
\bibitem{cssurvey} ``A comprehensive survey on client selection strategies in federated learning,'' \emph{Comput.\ Netw.}, 2024.
\bibitem{yin2018} D. Yin, Y. Chen, K. Ramchandran, and P. Bartlett, ``Byzantine-robust distributed learning: Towards optimal statistical rates,'' in \emph{Proc.\ ICML}, 2018.
\bibitem{krum} P. Blanchard, E. M. El Mhamdi, R. Guerraoui, and J. Stainer, ``Machine learning with adversaries: Byzantine tolerant gradient descent (Krum),'' in \emph{Proc.\ NeurIPS}, 2017.
\bibitem{ppo} J. Schulman, F. Wolski, P. Dhariwal, A. Radford, and O. Klimov,
``Proximal policy optimization algorithms,'' \emph{arXiv:1707.06347}, 2017.
\end{thebibliography}

\end{document}
